% Demo paper for the meertex website screenshots, after the reviews. Every
% name, number and reference in it is made up. ../main.tex is the version
% that was submitted.
\documentclass[acmsmall,review,anonymous]{acmart}

\settopmatter{printacmref=false}
\setcopyright{none}
\renewcommand\footnotetextcopyrightpermission[1]{}
\acmJournal{TOSEM}

\begin{document}

\title{Reproducible Energy Measurements for Software Engineering Research}

\author{Lena Vogt}
\affiliation{\institution{Example University}\country{Germany}}
\email{lena@example.com}

\begin{abstract}
Energy measurements of software are hard to repeat. Results change with the
machine, its temperature, the background load and the measurement tool. We
reviewed 134 studies published between 2015 and 2025 that report the energy
use of software. Only 18 of them report enough detail to repeat the
measurement. We propose a checklist of twelve items and show on four case
studies that it makes measurements repeatable within 3\,\%.
\end{abstract}

\maketitle

\section{Introduction}

Software engineering research increasingly reports the energy use of
programs, libraries and development practices. A result that cannot be
repeated on another machine says little about the software and a lot about
the setup it was measured on. Reviewers rarely ask for the setup, because
there is no common list of what it should contain.

\section{Literature Review}

We searched four digital libraries and two preprint servers for studies that
measure the energy of software and kept 134 of them. For each study we
recorded the hardware, the measurement tool, the number of repetitions and
whether idle power was subtracted. Only 18 studies report all four.

\section{The Checklist}

The checklist has twelve items in three groups. The first group describes
the machine, the second the measurement and the third the analysis. Each item
is a question a reader must be able to answer from the paper alone. We
tested the wording with eight researchers who had not seen it before.

\section{Case Studies}

We applied the checklist to four earlier studies and repeated their
measurements on a second machine. With the checklist the results agreed
within 3\,\%. Without it they differed by up to 40\,\%.

\section{Threats to Validity}

The case studies were picked by us and may be easier to repeat than the
average study. Two of the four use the same measurement tool.

\section{Conclusion}

Repeatable energy measurements need a description of the setup that is as
careful as the description of the software under test. The checklist is a
first step towards it.

\end{document}
